\section{Retranscription du besoin client}
\small

\subsection{Contexte}
La mairie de Saint-Mars-du-Désert est une collectivité territoriale dynamique et innovante, engagée dans une démarche de transition écologique et sociétale.
\newline

Dans ce cadre, la commune souhaite renforcer ses connaissances sur la faune locale en mettant en place un dispositif de surveillance et d'observation automatisée.
Ce projet s'inscrit dans une démarche plus large de constitution de l'Atlas de la Biodiversité Communale (ABC), outil essentiel pour inventorier, comprendre et valoriser la richesse écologique du territoire.
\newline

Le système envisagé aura pour objectif de capturer des photographies d’espèces animales présentes sur la commune. Dès leur acquisition, ces images seront transmises sans fil à un stockage distant afin d’être analysées par les services municipaux et les partenaires scientifiques et associatifs.
\newline

\subsection{Parties prenantes}
Dans le cadre de ce projet, l'entité ayant le rôle de \gloss{moa} sera la mairie de Saint-Mars-du-Désert.
L'interlocuteur ayant pour objectif de représenter la mairie dans ce projet sera M. Xavier LEPREVOST, élu à la mairie de Saint-Mars-du-Désert.
\newline

\begin{table}[h!]
   \centering
   \renewcommand{\arraystretch}{1.2} % espace entre les lignes
   \begin{tabular}{|c|c|c|c|}
      \hline
      \textbf{QUI}     & \textbf{RÔLE} & \textbf{MAIL}                         & \textbf{MOBILE}   \\
      \hline
      LEPREVOST Xavier & Client        & Xavier.leprevost@saintmarsdudesert.fr & +33 6 28 34 18 50 \\
      \hline
   \end{tabular}
   \caption{Informations des intervenants}
   \label{tab:intervenants}
\end{table}

\newpage
Un groupe de six étudiants occupera le rôle de \gloss{moe}.
Ce groupe est composé de trois étudiants ayant fait un BUT en électronique et informatique, un étudiant ayant fait une prépa PEIP (cursus préparatoire visant à préparer les élèves au cycle ingénieur Polytech), une étudiante ayant fait une licence en génie électrique et un étudiant ayant fait une licence en génie électrique, spécialité AII (automatique et informatique industrielle). \\[15pt]

\begin{table}[h!]
   \centering
   \renewcommand{\arraystretch}{1.2} % espace entre les lignes
   \begin{tabular}{|m{2.5cm}|m{5cm}|m{5cm}|c|}
      \hline
      \textbf{QUI}   & \textbf{RÔLE}                                & \textbf{MAIL}                       & \textbf{MOBILE}   \\
      \hline
      LE CALLET Ewen & Chef de projet                               & ewen.le-callet@saintmarsdudesert.fr & +33 6 95 15 93 30 \\
      \hline
      NOUMBA Jules   & Reponsable work package modules additionnels & jules.noumba@etu.univ-nantes.fr     & +33 6 11 83 82 77 \\
      \hline
      BUCHET Nolan   & Responsable work package hardware            & nolan.buchet@etu.univ-nantes.fr     & +33 6 83 68 16 46 \\
      \hline
      XU Xueying     & Responsable work package module caméra       & xueyig.xu@etu.univ-nantes.fr        & +33 7 82 67 76 72 \\
      \hline
      CLARK Davy     & Responsable work package stockage de données & davy.clark@etu.univ-nates.fr        & +33 7 85 78 48 46 \\
      \hline
      MENARD Joris   & Responsable workpakage transmission          & joris.menard@etu.univ-nantes.fr     & +33 6 71 61 61 96 \\
      \hline
   \end{tabular}
   \caption{Informations sur l'équipe de projet}
   \label{tab:intervenants}
\end{table}

\vspace{1.5em}
Le lien avec les partenaires scientifiques, associatifs et naturalistes sera assuré par la mairie, notamment par l’intermédiaire de Xavier Leprevost. Les données et photographies produites dans le cadre du projet seront mises à disposition de la mairie, qui pourra les transmettre aux partenaires concernés selon les besoins.